\documentclass[10pt,a4paper]{amsart}
\usepackage[margin=19mm]{geometry}
\usepackage{amsmath,amssymb,mathtools}
\usepackage[scr=rsfs,cal=euler]{mathalfa}
\usepackage{xfrac}
\usepackage[unicode,hidelinks,bookmarks=false]{hyperref}
\newcommand{\ie}{i.e.\ }
\newcommand{\cf}{cf.\ }
\newcommand{\resp}{respectively\ }
\newcommand{\Subsection}{\subsection}
\providecommand{\nobreakdash}{\nobreak\hbox{-}}
\setlength{\emergencystretch}{2em}
\multlinegap0pt
\usepackage{xcolor}
\usepackage{dsfont}
\usepackage{booktabs}
\usepackage{tikz}
\usepackage{natbib}
\usepackage{bm}
\usepackage{bbm}
\usepackage{amssymb}
\usepackage{csquotes}
\usepackage{cancel}
\usepackage{siunitx}
\usepackage{algorithm}\usepackage{algpseudocode}
\usepackage{caption}
\usepackage{enumerate}
\usepackage[normalem]{ulem}
\newtheorem{assumption}{Assumption}
\newtheorem{definition}{Definition}
\newtheorem{theorem}{Theorem}
\newtheorem{proposition}{Proposition}
\newtheorem{corollary}{Corollary}
\newcommand{\NN}{{\mathcal N \mathcal N}}
\newcommand{\R}{\mathbb{R}}
\newcommand{\calC}{{\mathcal C}}
\newcommand{\calK}{{\mathcal K}}
\newcommand{\calP}{{\mathcal P}}
\title{Operator Learning for Families of Finite-State Mean-Field Games}
\author{William Hofgard, Asaf Cohen, Mathieu Lauri\`ere}
\date{Source: arXiv:2602.13169v1 (CC BY 4.0)}
\begin{document}
\maketitle

\noindent\emph{Source and licence.} Operator Learning for Families of Finite-State Mean-Field Games. Version arXiv:2602.13169v1 (CC BY 4.0), distributed by arXiv under the Creative Commons Attribution 4.0 International licence, http://creativecommons.org/licenses/by/4.0/, as recorded in the arXiv licence metadata for this exact version and shown on the abstract page https://arxiv.org/abs/2602.13169v1. Full licence notice: \texttt{licenses/arxiv-2602.13169-CC-BY-4.0.txt}.\par\smallskip

\noindent\textit{Editorial scope.} Counted unit: the article's proposition prop:jiao-approx-bound (source0.tex lines 469--475). The article's complete original proof of it is delivered with it (source0.tex lines 1084--1094, one \texttt{proof} environment of 11 lines, which the article places away from the statement). No mathematics is written, repaired or re-derived: the statement and the proof are the article's own lines, in the article's own notation, and no retained line is substituted or re-flowed.  Retained uncounted as the source-local context the proof needs: lines 338--345; lines 357--364; lines 386--389; lines 394--401; lines 454--461; lines 477--484.  Nothing was omitted from any retained span, so the package carries no omission record.  No retained line contains a citation, so the wrapper carries no bibliography and no bibliography entry.  No retained line contains a figure environment, an image-inclusion command or drawing code, so no image is delivered and the package declares no asset dependency.  Editorial formatting only, in addition: an AMS \texttt{amsart} preamble that restates the article's own declarations and macros used by the retained lines, namely the environments \texttt{assumption}, \texttt{definition}, \texttt{theorem}, \texttt{proposition}, \texttt{corollary} on separate counters, together with the standard packages those lines need.  The retained environments print in this document as Assumption 1, Definition 1, Theorem 1, Proposition 1, Corollary 1 in order of appearance, whereas the article prints the same items with the numbering of its own full document.  The complete native source of the article as distributed in the arXiv e-print for this exact version is bundled unchanged as \texttt{sources/194622/source0.tex}.
\begin{align}
\label{eqn:mfg}
\begin{cases}
    \frac{d}{dt} u (t,x) + \bar H(x, \mu(t), \Delta_x u (t,\cdot)) = 0, & (t,x) \in [0, T] \times [d] \qquad \textbf{(HJB)} \\
    \frac{d}{dt} \mu (t,x) = \sum_{y\in [d]} \mu (t,y) \gamma^*_x(y,\Delta_y u (t,\cdot)), & (t,x) \in [0, T] \times [d] \qquad \textbf{(KFP)} \\
    \mu (0, x) = \eta(x), \quad u(T, x) = g(x, \mu(T)), & x \in [d],
\end{cases}
\end{align}

\begin{assumption}
\label{asmp:mfg-uniqueness}
    The minimizer $\gamma^*(x, p)$ of the Hamiltonian $H$ is unique. Moreover, $H$ is strictly concave in $p$ and twice continuously differentiable with Lipschitz second derivatives. Finally, the costs $F$ and $g$ are continuously differentiable with Lipschitz derivatives, and both are \textbf{Lasry--Lions monotone} in the sense that for both $\phi = F, g$,
    \begin{align}\label{eqn:LL-monotone}
    \textstyle
        \sum_{x \in [d]} (\phi(x, \eta) - \phi(x, \hat{\eta})) (\eta_x - \hat{\eta}_x) \geq 0, \qquad \eta, \hat{\eta} \in \mathcal{P}([d]).
    \end{align} 
\end{assumption}

\begin{definition}
\label{def:flow-map}
Given a parametrized family of terminal conditions, the \textbf{flow map} $\Phi : [0, T] \times \calP([d]) \times \calK \to \R^d$ is defined by $\Phi(t, \eta, \kappa) := u^{\eta, \kappa}(t)$, where $u^{\eta, \kappa}$ is the value function for the MFG system~\eqref{eqn:mfg} with initial distribution $\eta \in \calP([d])$ and terminal cost $g_\kappa$.
\end{definition}

\begin{assumption}
\label{asmp:param-terminal-cost}
There exists a compact set of parameters $\calK \subseteq \R^k$ such that for all $\kappa \in \calK$, the $g_\kappa : [d] \times \calP([d]) \to \R$ satisfies Assumption~\ref{asmp:mfg-uniqueness}. Moreover, for any $\kappa, \kappa' \in \calK$, there exists a constant $C > 0$ such that
$
    |g_{\kappa}(x, \mu) - g_{\kappa'}(x, \mu)| \leq C|\kappa - \kappa'|,
$
uniformly in $(x, \mu) \in [d] \times \calP([d])$.
\end{assumption}

\begin{theorem}
\label{thm:flow-time-regularity}
Under Assumptions~\ref{asmp:mfg-uniqueness} and ~\ref{asmp:param-terminal-cost}, the flow map $\Phi : [0, T] \times \calP([d]) \times \calK \to \R^d$, given by $\Phi(t, \eta, \kappa) = u^{t_0, \eta, \kappa}(t, \cdot)$, is jointly Lipschitz in its inputs: there exists $C > 0$ such that
\begin{align*}
    |\Phi(t, \eta_1, \kappa_1) - \Phi(s, \eta_2, \kappa_2)| \leq C(|t - s| + |\eta_1 - \eta_2| + |\kappa_1 - \kappa_2|)
\end{align*}
for all $(t, \eta_1, \kappa_1), (s, \eta_2, \kappa_2) \in [0, T] \times \calP([d]) \times \calK$.
\end{theorem}

\begin{corollary}
\label{cor:appx-flow-map}
Assume that Assumptions~\ref{asmp:mfg-uniqueness} and~\ref{asmp:param-terminal-cost} hold. Then, for any $K \geq 1$ and $\varepsilon > 0$, there exists a neural network $\Psi \in \NN(W, L, K)$ with weight bound $K$, width $W \geq c(\mathrm{diam}(\calK), T) d K^{(2(d + k) + 3)/(2(d + k) + 4)}$, and depth $L \geq 2 \lceil \log(d + k + 1)\rceil + 2$ such that
\begin{align*}
    \|\Phi - \Psi\|_{\calC([0, T] \times \calP([d]) \times \calK)} \leq C(\mathrm{diam}(\calK), T) K^{-1/(d + k + 2)} + \varepsilon,
\end{align*}
where $\Phi : [0, T] \times \calP([d]) \times \calK \to \R^d$ is the flow map from Definition~\ref{def:flow-map}.
\end{corollary}

\begin{proposition}
\label{prop:jiao-approx-bound}
There exists constants $c, C > 0$ such that for any $K \geq 1$, $W \geq c K^{(2d+1) / (2d+2)}$, and $L \geq 2 \lceil \log(d)\rceil + 2$, the worst-case approximation error of the class $\NN(W, L, K)$ for $\Phi \in \calC^{0, 1}([0,1]^d)$ satisfies:
$
    \sup_{\Phi \in \calC^{0, 1}([0,1]^d)} \inf_{\Psi \in \NN(W, L, K)} \|\Phi - \Psi\|_{\calC([0,1]^d)} \leq C K^{-1/(d + 1)}.
$
\end{proposition}

\begin{proof}[Proof of Corollary~\ref{cor:appx-flow-map}]
First, to extend from the setting of scalar regression, as is the case in Proposition~\ref{prop:jiao-approx-bound}, to vector-valued regression, we note that if a scalar Lipschitz function can be uniformly approximated up to an error $\varepsilon > 0$ by a network with weight bound $K$ and width $W$, then an $\R^d$ valued Lipschitz function $\Phi$ can be uniformly approximated by a network with weight bound $K$ and width $d W$. Indeed, we can simply approximate each coordinate of $\Phi$ by a network of width $W$ and stack the resulting networks to obtain the desired approximator, which will have width $d W$.

Now, by Theorem~\ref{thm:flow-time-regularity}, the flow map $\Phi : [0, T] \times \calP([d]) \times \calK \to \R^d$ belongs to $\calC^{0,1}([0, T] \times \calP([d])\times \calK)$. From this, we can apply Proposition~\ref{prop:jiao-approx-bound} directly upon scaling the domain $[0, T] \times \calP([d])\times \calK$ to lie entirely within the $(d + k +1)$-dimensional unit cube.

To carry out this scaling, we embed $\calP([d]) \hookrightarrow [0, 1]^d$, scale $\calK$ to lie in the set $[0,1]^k$, and scale the interval $[0, T]$ to lie in the interval $[0, 1]$. The natural embedding $\calP([d]) \hookrightarrow [0, 1]^d$ is simply given by viewing
\begin{align*}
    \calP([d]) = \left\{\eta \in \R^d : \sum_{i=1}^d \eta_i = 1, \quad \eta_i \geq 0 \text{ for all } i = 1, \ldots, d\right\}.
\end{align*}
This rescaling may incur constants that depend on the diameter of $\calK$, denoted by $\mathrm{diam}(\calK)$, and the final time $T$. Importantly, it is always possible for finite $T > 0$ and compact $\calK \subset \R^k$. The result then follows upon applying Proposition~\ref{prop:jiao-approx-bound}, replacing $d$ with $d + k + 1$ therein. As noted above, the universal constants $c, C > 0$ obtained in Proposition~\ref{prop:jiao-approx-bound} must also be replaced by constants $c(\mathrm{diam}(\calK), T), C(\mathrm{diam}(\calK), T) > 0$ that depends on $\calK$ and $T$.
\end{proof}

\end{document}
